%%%%%%%%%%%%%%%%%%%%%%% file template.tex %%%%%%%%%%%%%%%%%%%%%%%%%
%
% This is a general template file for the LaTeX package SVJour3
% for Springer journals.          Springer Heidelberg 2010/09/16
%
% Copy it to a new file with a new name and use it as the basis
% for your article. Delete % signs as needed.
%
% This template includes a few options for different layouts and
% content for various journals. Please consult a previous issue of
% your journal as needed.
%
%%%%%%%%%%%%%%%%%%%%%%%%%%%%%%%%%%%%%%%%%%%%%%%%%%%%%%%%%%%%%%%%%%%
%
% First comes an example EPS file -- just ignore it and
% proceed on the \documentclass line
% your LaTeX will extract the file if required
% \begin{filecontents*}{example.eps}
%!PS-Adobe-3.0 EPSF-3.0
%%BoundingBox: 19 19 221 221
%%CreationDate: Mon Sep 29 1997
%%Creator: programmed by hand (JK)
%%EndComments
% gsave
% newpath
%   20 20 moveto
%   20 220 lineto
%   220 220 lineto
%   220 20 lineto
% closepath
% 2 setlinewidth
% gsave
%   .4 setgray fill
% grestore
% stroke
% grestore
% \end{filecontents*}
% \RequirePackage{fix-cm}
% \documentclass[smallextended,referee,envcountsect]{svjour3}       % onecolumn (second format)
\documentclass[smallextended,envcountsect]{svjour3}       % onecolumn (second format)
\smartqed  % flush right qed marks, e.g. at end of proof
\usepackage{graphicx}
% \usepackage{textcomp}
% \usepackage{subfigure}
\usepackage[center]{caption}
\usepackage{amsfonts}
% \usepackage{pgfplots}
\usepackage{algorithm}
\usepackage[noend]{algorithmic}
\usepackage{amsmath}

\newtheorem{assumption}{Assumption}
\newcommand{\len}[1]{\text{len}{(#1)}}

\journalname{To be decided}

\begin{document}

\title{Confluent flow and other approximations to the Pooling Problem}
\titlerunning{Confluent flow and other approximations to the Pooling Problem}

\author{Dag Haugland \and Eligius M.T. Hendrix \and \'{I}sak P\'{a}ll Gestsson}

\institute{Dag Haugland, Corresponding author \at 
Department of Informatics, University of Bergen, Norway \email{dag.haugland@ii.uib.no}
\and
Eligius M.T. Hendrix \at
              Department of Computer Architecture, Universidad de M\'alaga and Operations Research and Logistics Group, Wageningen University \email{eligius.hendrix@wur.nl}\and \'{I}sak P\'{a}ll Gestsson \at Gassco}

\date{Received: date / Accepted: date}

\maketitle

\begin{abstract}
About the computational challenge of the pooling problem and finding approximate solutions by means of confluent or influent flow
\end{abstract}
\keywords{Pooling problem \and complexity \and global optimization}
\subclass{90C26 \and 90C35 \and 90C39 \and 68Q25}


%-------------------------------------------------------------------------------
\section{Introduction}
\label{intro}

Taken from an earlier paper, to be revised:
Traditional network flow models, like the minimum cost network flow problem, rely on considerations of the flow as a homogeneous commodity.
Variations in composition is however relevant in many industrial and logistics applications of the models,
including petroleum refining and pipeline transportation of natural gas.
Crude oils supplied to a refinery have contamination levels depending on their sources of origin.
Major components of natural gas are methane, ethane, butane and propane, but the proportions in which they occur are not equal for all gas wells.
For environmental or technical reasons, requirements to the final composition are often imposed at the sink nodes of the flow network.
Whenever constraints of this kind apply, it is therefore essential for network flow models to recognize not only the total product flow,
but also how the flow composition evolves from network sources to sinks.
In particular, the models must reflect updates of the composition at nodes where differently composed flows are pooled.
This provides a computationally challenging problem referred to as the \textsc{Pooling Problem}.

The \textsc{Pooling Problem} exhibits strong resemblance with well-studied logistics models like the \textsc{Minimum Cost Flow Problem} and the
\textsc{Transportation Problem}.
Another classical related problem is Stigler's \textsc{Diet Problem} \cite{stigler45,dantzig90},
asking for the minimum cost diet requiring a certain amount of specified nutrients.
The problem is easy to extend to a multi-diet version where groups of people have each their own requirements.
Nutritional contents differ between sources, that is, foods of which diets can be composed,
and the optimal mix must be found for each sink (group of people).
Conceptually, the multi-diet version coincides with the \textsc{Pooling Problem} in bipartite networks.
Thanks to this simple network structure, with exclusively direct links between sources and sinks, diet problems are nothing but linear programs.
However, when extending the problem to more general network structures, this is no longer the case.
Absence of linear formulations has immediate computational consequences.
While minimum cost flow instances where the network size reaches a six digit number of arcs and nodes can be solved within hours, or even minutes,
exact solution of \textsc{Pooling Problem} instances with hundreds of nodes appears to be unrealistic.

The \textsc{Pooling Problem} has been recognized for decades as a challenge where linear programming approaches fail.
Bounds on the flow composition frequently induce bilinear constraints in otherwise linear models,
and require corresponding solution methods.
An approach based on iterated solution of linearized model versions has gained popularity especially in the oil refining industry,
and has been analyzed experimentally by Haverly \cite{haverly78}.
A major conclusion is that methods of this kind may produce suboptimal solutions.
Such observations helped establishing a consensus, already in early days of research in the \textsc{Pooling Problem}, that global optimality is achievable only in small instances.

Owing to the anticipated intractability of the problem, early research was mainly directed towards heuristic methods \cite{haverly78,baker85,dewitt89},
where iterative linearization is the core idea.
Floudas and Visweswaran \cite{floudas90} reported the first exact solution algorithm for the \textsc{Pooling Problem}.
Later, algorithms based on branch-and-bound \cite{audet04,foulds92,misener11,sahinidis05,visweswaran96},
Lagrangian relaxation \cite{adhya99,almutairi09,bental94}, and integer programming \cite{dey15} have been studied.
The industrial relevance of the problem, notably in petroleum refining \cite{baker85,dewitt89},
the food industry \cite{kallrath05}, and in waste water processing \cite{galan98,misener10},
has been acknowledged by many authors.
The \textsc{Pooling Problem} survey by Misener and Floudas \cite{misener09} has a comprehensive list of references to work in this area.

The \textsc{Standard Pooling Problem} (SPP) is defined in terms of a networks with three node layers, referred to as sources, pools and sinks (sinks), respectively.
Each arc links either a source to a pool or a pool to a sink.
Quality constraints in terms of bounds on the composition are imposed at the sinks.
Each bound relates to the relative content of a flow component, referred to as the component.
After NP-hardness of the SPP was formally proved \cite{alfaki13}, Dey and Gupte \cite{dey15} showed that the problem is as hard to approximate as the maximum clique problem.
The problem remains NP-hard for networks with only two sources and two sinks,
and only one flow component subject to constraints \cite{haugland16}.

The remainder of the text is organized as follows:
Mathematical notation and formulation of the \textsc{Pooling Problem} are provided in Section \ref{sec:problem}.

To be written.

%-----------------------------------------------------------------------
\section{Mathematical Notation and Problem Formulation}\label{sec:problem}

To be revised:

Let $D=(N,A)$ be a directed acyclic graph,
where the node set $N=S\cup P\cup T$ consists of the disjoint subsets $S$ of \emph{sources}, $P$ of \emph{pools}, and $T$ of \emph{sinks}.
Let $K$ be a finite set, and refer to its elements as \emph{flow components} or simply \emph{components}.
In the sequel, elements of vectors and matrices defined over the sets $N$ and $K$
are identified by a subscript if they correspond to $N$, and by a superscript if they correspond to $K$.
Complexity analyses in later sections require that network parameters are defined as rational rather than real numbers.
Define the vectors of unit arc costs $\boldsymbol{c}\in\mathbb{Q}^{A}$ and arc capacities $\boldsymbol{h}\in\mathbb{Q}_+^{A}$.
For integers $j$ and $n$,
let $\boldsymbol{j}^n$ denote the vector with all its $n$ entries equal to $j$,
and let $\boldsymbol{I}^n$ denote the $(n\times n)$-identity matrix.

At all sources $s\in S$, flow with a given concentration of each component $k\in K$ may enter $D$, and flow leaving the network at a sink $t\in T$ must satisfy upper concentration bounds.
Define the matrices of \emph{concentration} $\boldsymbol{Q}\in\mathbb{Q}^{K\times S}$ and \emph{concentration bounds}
$\boldsymbol{U}\in\mathbb{Q}^{K\times T}$.
For $s\in S$ and $t\in T$, let $\boldsymbol{q}_s=\left(q_s^k\right)_{k\in K}$ and $\boldsymbol{u}_t=\left(u_t^k\right)_{k\in K}$ be the columns of $\boldsymbol{Q}$ and $\boldsymbol{U}$, respectively.
That is, $q_s^k$ is the concentration of component $k$ of the flow entering $D$ at $s$,
and $u_t^k$ is the upper bound on the concentration of the flow leaving $D$ at $t$.

For each pool $p\in P$, the set of sources from which there exists a path to $p$ in $D$ is denoted $S_p$.
Analogously, $T_p$ is the set of sinks to which there is a path from $p$ in $D$.
The sets of upstream and downstream neighbors of node $j\in N$ are denoted $N_j^-$ and $N_j^-$, respectively.
Without loss of relevant generality, it is assumed that $N_s^-=N_t^+=\emptyset$ for every $s\in S$ and $t\in T$.
The sets of pools adjacent to source $s$ and sink $t$ are denoted
$P_s=P\cap N_s^+$ and $P_t=P\cap N_t^-$, respectively.
% The arc set $\mathcal{A}$ has a partition $\left(\mathcal{A}_S,\mathcal{A}_T\right)$,
% where $\mathcal{A}_S=\mathcal{A}\cap\left(S\times P\right)$ and $\mathcal{A}_T=\mathcal{A}\cap\left(P\times T\right)$
% denote the sets of source-incident and sink-incident arcs, respectively.

A vector $\boldsymbol{f}\in\mathbb{R}_+^{A}$ is called a \emph{flow} in $D$ iff it satisfies the capacity and flow conservation constraints:
\begin{eqnarray}
           & f_{ij}\leq h_{ij} & (i,j)\in A, \label{eq:cap} \\
           & \sum_{j\in N_p^-}f_{jp}-\sum_{j\in N_p^+}f_{pj} = 0~~~~ & p\in P. \label{eq:flowcons}
\end{eqnarray}

\noindent
We refer to the set $F(D,\boldsymbol{h}) = \left\{\boldsymbol{f}\in\mathbb{R}_+^{A}: \mbox{(\ref{eq:cap})--(\ref{eq:flowcons}) are satisfied}\right\}$
as the \emph{flow polytope} of $D$.
For each $\boldsymbol{f}\in F(D,\boldsymbol{h})$,
define $\boldsymbol{W}=\left(\boldsymbol{w}_j\right)_{j\in N}\in\mathbb{R}^{K\times N}$ to be a
\emph{concentration matrix} of $\boldsymbol{f}$ iff $\boldsymbol{w}_s=\boldsymbol{q}_s$ ($s\in S$) and
\begin{equation} \label{eq:totqual}
  \boldsymbol{w}_p\sum_{j\in N_p^-}f_{jp}=\sum_{j\in N_p^-}\boldsymbol{w}_jf_{jp},\ p\in P\cup T.
\end{equation}
That is, the concentration $\boldsymbol{w}_p=\left(w_p^k\right)_{k\in K}$ of component $k$ in the flow blended at a pool/sink $i$,
is a convex combination of the entering concentrations $\boldsymbol{w}_j$ ($j\in N_i^-$),
with convexity weights $f_{ji}/\sum_{j'\in N_i^-}f_{j'p}$.
Condition \eqref{eq:totqual} on $\boldsymbol{w}$ is referred to as the \emph{linear blending} property.
At sinks $t\in T$ where $\sum_{i\in N_t^-}f_{it}>0$, the upper concentration bounds $u_t^k$ apply,
and linear blending yields $\boldsymbol{w}_t=\sum_{j\in N_t^-}\boldsymbol{w}_jf_{jt}/\sum_{j\in N_t^-}f_{jt}\leq \boldsymbol{u}_t$.
Thus, $\boldsymbol{f}\in F\left(D,\boldsymbol{h}\right)$ is said to be a \emph{concentration-feasible} flow in
$D$ iff $\sum_{j\in N_t^-}\boldsymbol{w}_jf_{jt}\leq \boldsymbol{u}_t\sum_{j\in N_t^-}f_{jt}$ ($t\in T$)
for any concentration matrix $\boldsymbol{W}$ of $\boldsymbol{f}$.

\begin{problem} \label{prob:pp}
The \textsc{Pooling Problem}: Find a concentration-feasible flow $\boldsymbol{f}$ in $D$ minimizing the total costs $\boldsymbol{c}\cdot\boldsymbol{f}$.
\qed
\end{problem}

%-----------------------------------------------------------------
\section{Cofluent flow and extensions} \label{sec:conf}
In the \textsc{Minimum Cost Flow Problem} (MCFP), the flow leaving a source and the flow entering a sink are subject to bounds.
Let the bounds in question, referred to as the \emph{supply} at source $s\in S$ and the \emph{demand} at sink $t\in T$, respectively, be denoted $b_s$ and $b_t$,
and let $\boldsymbol{b}=\left(b_u\right)_{u\in S\cup T}$.
Traditionally in MCFP, the supply represents an upper bound, whereas demand is a \emph{lower} bound on the sink flow.
This contrasts SPP, in which both bounds are upper, with the implication that zero flow is feasible in all instances.
In the current work, we will consider both these options, and state explictly whether demand should be considered an upper or a lower bound.

Let
\[
F^{\geq}(D,\boldsymbol{h},\boldsymbol{b}) = \left\{\boldsymbol{f}\in F(D,\boldsymbol{h}): \sum_{j\in N_s^+}f_{sj}\leq b_s (s\in S), \sum_{j\in N_t^-}f_{jt}\geq b_t (t\in T)\right\}
\]
and
\[
F^{\leq}(D,\boldsymbol{h},\boldsymbol{b}) = \left\{\boldsymbol{f}\in F(D,\boldsymbol{h}): \sum_{j\in N_s^+}f_{sj}\leq b_s (s\in S), \sum_{j\in N_t^-}f_{jt}\leq b_t (t\in T)\right\}
\]
be the flow polytopes corresponding to the two variants of the MCFP.
In the current article, the notion of \emph{confluent} flow \cite{chen06} and related concepts are extensively exploited:
\begin{definition} \label{def:conf}
A flow $\boldsymbol{f}\in F(D,\boldsymbol{h})$ is said to be \emph{confluent} if, at all nodes $i\in S\cup P$, $f_{ij}>0$ for at most one $j\in N_i^+$.
\qed
\end{definition}

That is, confluence means that while flow streams are allowed to merge at pools, it has to leave the pool along a unique arc.
Also all flow leaving a source must be allocated to the same arc.
The inflow counterpart of confluent flow, where merging is disallowed, is defined analogously:
\begin{definition} \label{def:inf}
A flow $\boldsymbol{f}\in F(D,\boldsymbol{h})$ is said to be \emph{influent} if, at all nodes $i\in P\cup T$, $f_{ji}>0$ for at most one $j\in N_i^-$.
\qed
\end{definition}

The two concepts are combined disjunctively:
\begin{definition} \label{def:confinf}
A flow $\boldsymbol{f}\in F(D,\boldsymbol{h})$ is said to be \emph{confluent-influent} if, at all nodes $p\in P$, either $f_{pj}>0$ for at most one $j\in N_p^+$
or $f_{jp}>0$ for at most one $j\in N_p^-$.
\qed
\end{definition}

Observe that the condition in Definition~\ref{def:confinf} applies to pools only, since no flow enters sources or leaves sinks.

\begin{problem} \label{prob:conf}
The \textsc{Minimum Cost Confluent Flow Problem} (MCCFP):
Find a flow $\boldsymbol{f}\in F^{\geq}(D,\boldsymbol{h},\boldsymbol{b})$ satisfying Definition~\ref{def:conf} such that the cost $\boldsymbol{c}\cdot\boldsymbol{f}$ is minimized.
\qed
\end{problem}
\noindent
Analogously, in correspondence with Definitions~\ref{def:inf} and \ref{def:confinf}, respectively, the 
\textsc{Minimum Cost Influent Flow Problem} (MCIFP) and the \textsc{Minimum Cost Confluent-Influent Flow Problem} (MCCIFP) are defined.

\subsection{Complexity of fluency-constrained extensions of MCFP} \label{sec:compl}
A modified version of MCCFP, in which the largest arc flow is to be minimized, is known to be NP-hard \cite{chen06}, and even NP-hard to approximate within a factor
$\tfrac{4}{3}$.
By a simple reduction from the \textsc{Subset Sum Problem},
it is next shown that MCCFP as introduced in Definition~\ref{prob:conf} is NP-hard, even in a much restricted instance class.
\begin{proposition} \label{prop:conf}
MCCFP is NP-hard in instances where $D$ is bipartite, $P=\emptyset$, and $|T|=2$.
\end{proposition}
\begin{proof}
Consider an instance $(b_1,\ldots,b_n,d)$ of the \textsc{Subset Sum Problem}, with positive integers $1\leq b_1\leq\cdots\leq b_n\leq d$.
Let $D$ be the complete bipartite digraph with node sets $S=\{s_1,\ldots,s_n\}$ and $T=\{t_0,t_1\}$.
Let the supply at $s_i$ be $b_i$ ($i=1,\ldots,n$), while the demand at $t_0$ and $t_1$ is, respectively, $\sum_{i=1}^nb_i-d$ and $d$.
Both arcs leaving $s_i$ have capacity $b_i$ and arbitrary cost.
Because total supply equals total demand, $\boldsymbol{f}\in F^{\geq}(D,\boldsymbol{h},\boldsymbol{b})$ only if $b_i$ flow units leave source $s_i$,
$d$ units enter sink $t_1$, and $\sum_{i=1}^nb_i-d$ units enter sink $t_0$.
Confluence implies that all the flow leaving $s_i$ enters exactly one of the sinks $t_0$ or $t_1$.
Thus, the MCCFP instance $(D,\boldsymbol{h},\boldsymbol{b},\boldsymbol{c})$ is feasible if and only if there exists a set $S_1\subseteq\{1,\ldots,n\}$ such that $\sum_{i\in S_1}b_i=d$.
\qed
\end{proof}

NP-hardness in restricted instance classes also applies to the influent and confluent-influent counterparts of MCCFP:
\begin{proposition} \label{prop:inf}
MCIFP is NP-hard in instances where $D$ is bipartite, $P=\emptyset$, and $|S|=2$.
\end{proposition}
\begin{proof}
Consider the instance of the \textsc{Subset Sum Problem} introduced in the proof of Proposition~\ref{prop:conf}.
Let $D$ be the complete bipartite digraph with node sets $S=\{s_0,s_1\}$ and $T=\{t_1,\ldots,t_n\}$.
Let the demand at $t_i$ be $b_i$ ($i=1,\ldots,n$), while the supply at $s_0$ and $s_1$ is, respectively, $\sum_{i=1}^nb_i-d$ and $d$.
Finally, $c_{st_i}$ is arbitrary and $h_{st_i}=b_i$ for $s\in\{s_0,s_1\}$.
Following the arguments in the proof of Proposition~\ref{prop:conf}, $(D,\boldsymbol{h},\boldsymbol{b},\boldsymbol{c})$ is a feasible MCIFP instance
if and only if $\sum_{i\in S_1}b_i=d$ for some $S_1\subseteq\{1,\ldots,n\}$.
\qed
\end{proof}

Proposition~\ref{prop:conf} also applies if MCCFP is defined with the demand as upper bounds.
That is, if $F^{\geq}$ is replaced by $F^{\leq}$ in Problem~\ref{prob:conf}.
Consider the digraph introduced in the proof of the proposition, and let the cost $c_{st}=-2$ for all $s=s_1,\ldots,s_n$ and $t=t_1$.
For $t=t_0$, let $c_{st}=-1$.
Hence, the minimum cost is $-t-\sum_{i=1}^nb_i$ if and only if $\sum_{i\in S_1}b_i=d$ for some $S_1\subseteq\{1,\ldots,n\}$.
Analogously, Proposition~\ref{prop:inf} applies even if MCIFP is defined with respect to $F^{\leq}$ rather than $F^{\geq}$.

Let $\len{D}$ denote the maximum length of paths from a source to a sink in $D$.
In bipartite digraphs consisting only of sources and sinks (i.e, $\len{D}=1$), any flow is confluent-influent.
Introducing an intermediate set $P$ of nodes renders MCCIFP NP-hard, even if there are no arcs between any two nodes in $P$.
\begin{proposition} \label{prop:confinf}
MCCIFP is NP-hard in instances where $\len{D}=|T|=2$.
\end{proposition}
\begin{proof}
Consider the instance of the \textsc{Subset Sum Problem} introduced in the proof of Proposition~\ref{prop:conf}.
Let $D$ be the complete bipartite digraph with node sets $S=\{s_1,\ldots,s_n,s_1',\ldots,s_n'\}$, $P=\{p_1,\ldots,p_n\}$, and $T=\{t_0,t_1\}$.
All arcs of $D$ are incident to $P$, where $N_{p_i}^-=\{s_i,s_i'\}$ and  $N_{p_i}^+=\{t_0,t_1\}$.
% Node $p_i\in P$ has entering arcs $(s_i,p_i)$ and $(s_i',p_i)$, and leaving arcs $(p_i,t_0)$
Let the supply at $s_i$ and $s_i'$ be $b_i$ ($i=1,\ldots,n$), while the demand at $t_0$ and $t_1$ is, respectively, $2\sum_{i=1}^nb_i-2d$ and $2d$.
Both arcs entring $p_i$ have capacity $b_i$, whereas the leaving arcs have capacity $2b_i$. Arc costs are arbitrary.

Consider a feasible flow $\boldsymbol{f}$ in the MCCIFP instance $(D,\boldsymbol{h},\boldsymbol{b},\boldsymbol{c})$ hence defined.
Then, since total supply equals both the total capacity of arcs leaving $S$ and total demand, $f_{sp_i}=h_{sp_i}>0$ for all $s=s_i,s_i'$ and $i=1,\ldots,n$.
It follows that $\boldsymbol{f}$ is confluent, as it is not influent at any $p_i$.
Hence, all $2b_i$ units of flow leaving $p_i$ are sent to the same sink, $t_0$ or $t_1$.
Since $t_1$ receives $2d$ units of flow, there is a set $I_1\subseteq \{1,\ldots,n\}$ such that $\sum_{i\in I_1}f_{p_it_1}=2\sum_{i\in I_1}{b_i}=2d$.

Conversely, if $\sum_{i\in I_1}b_i=d$ for some $I_1\subseteq\{1,\ldots,n\}$, then the flow $\boldsymbol{f}$, where $f_{s_ip_i}=f_{s_i'p_i}=b_i$ ($i=1,\ldots,n\}$),
$f_{p_it_0}=\left\{\begin{array}{lr}2b_i, & i\notin I_1 \\ 0, & i\in I_1,\end{array}\right.$ and
$f_{p_it_1}=\left\{\begin{array}{lr}2b_i, & i\in I_1 \\ 0, & i\notin I_1,\end{array}\right.$ is feasible in the MCCIFP instance.
\qed
\end{proof}

\subsection{Weak fluency conditions} \label{sec:weak}

Different digraphs may be eqivalent in the perspective of the MCFP.
Splitting any node into two new ones joined by an arc with sufficently large capacity and zero cost, renders the MCFP instance unchanged.
In contrast, the operation largely alters the MCCIFP instance, since any flow in the modified network is confluent-influent.
A concept that is more resilient to likely graph alterations is introduced next.

For any flow $\boldsymbol{f}$ in $D$, let $D(\boldsymbol{f})=\left(N(\boldsymbol{f}), A(\boldsymbol{f})\right)$ denote the subgraph of $D$,
where $A(\boldsymbol{f})\subseteq A$ consists of arcs $(i,j)\in A$ where $f_{ij}>0$, and
$N(\boldsymbol{f})\subseteq N$ consists of nodes in $N$ with at least one incident arc in $A(\boldsymbol{f})$.

\begin{definition} \label{def:weakconfinf}
A flow $\boldsymbol{f}$ in $D$ is said to be \emph{weakly confluent} if, for all $i\in N(\boldsymbol{f})$,
$D(\boldsymbol{f})$ has paths from $i$ to at most one sink $t\in T\cap N(\boldsymbol{f})$.
If $D(\boldsymbol{f})$ has paths to $i$ from at most one source $s\in S\cap N(\boldsymbol{f})$, then $\boldsymbol{f}$ is said to be \emph{weakly influent}.
If either of the conditions holds at all $i\in N(\boldsymbol{f})$, then $\boldsymbol{f}$ is \emph{weakly confluent-influent}.
\qed
\end{definition}

It is readily observed that if $\boldsymbol{f}$ is confluent, then it is also weakly confluent.
Analogous claims hold for the other two fluency conditions.
Further, in bipartite digraphs, weak confluence (weak influence) implies confluence (influence).
Finally, if $\len{D}=2$ (see the proof of Proposition~\ref{prop:confinf}), then a weakly confluent-influent flow in $D$ is also confluent-influent.
For these reasons, Propositions~\ref{prop:conf}-\ref{prop:confinf} apply even if the fluency condition is to be understood in the weak sense.

\subsection{Single-commodity flow formulations for weakly confluent-influent flow} \label{sec:sciMILP}
Combining the single-commodity flow variables $\boldsymbol{f}$ with binary variables gives a natural formulation of the weak MCCIFP.
For all $p\in P$, $s\in S_p$, and $t\in T_p$, let $x_p^s=1$ signify that all paths from $S$ to $p$ in $D(\boldsymbol{f})$ start at $s$.
Analogously, $x_p^t=1$ means that all paths from $p$ to $T$ terminate at $t$.
For notational convenience, define the constants $x_u^v=\left\{\begin{array}{lr}1, & u=v \\ 0, & u\neq v,\end{array}\right.$r
for all $u,v\in S\cup T$, and $x_p^u=0$ for all $p\in P$ and $u\in (S\setminus S_p)\cup(T\setminus T_p)$.

Weakly confluent-influent flow dictates $\sum_{s\in S_p}x_p^s+\sum_{t\in T_p}x_p^t=1$ at all $p\in P$.
Consistency between $\boldsymbol{f}$ and $\boldsymbol{x}=\left(x_p^u\right)_{p\in P, u\in S_p\cup T_p}$ implies that for all arcs $(i,j)\in A(\boldsymbol{f})$
and all $s\in S_i$, it holds that $x_i^s=1$ (node $i$ receives flow only from source $s$) if $x_j^s=1$
(node $j$ receives flow only from source $s$).
Hence, $f_{ij}=0$ if $x_i^s=0$ and $x_j^s=1$, which is imposed by $f_{ij}\leq h_{ij}\left(1-x_j^s+x_i^s\right)$.
A symmetric argument suggests the constraint $f_{ij}\leq h_{ij}\left(1-x_i^t+x_j^t\right)$ for each $t\in T_j$.
A basic MILP model thus reads:
\begin{eqnarray}
\min & \sum_{(i,j)\in A}c_{ij}f_{ij} \label{eq:objMCCIFP} \\
     & \sum_{s\in S_p}x_p^s+\sum_{t\in T_p}x_p^t = 1 & p\in P \label{eq:convMCCIFP} \\
     & f_{ij}\leq h_{ij}\left(1-x_j^s+x_i^s\right) & (i,j)\in A, s\in S_j \label{eq:sMCCIFP} \\
     & f_{ij}\leq h_{ij}\left(1-x_i^t+x_j^t\right) & (i,j)\in A, t\in T_i \label{eq:tMCCIFP} \\
     & \boldsymbol{f}\in F^{\geq}(D,\boldsymbol{h},\boldsymbol{b}) \label{eq:fMCCIFP} \\
     & x_p^u\in\{0,1\} & p\in P, u\in S_p\cup T_p \label{eq:xMCCIFP}
\end{eqnarray}

For arcs $(i,j)\in A$, where $i\in S$, \eqref{eq:sMCCIFP} is simplified to $f_{ij}\leq h_{ij}(1-x_j^s)$ for all $s\neq i$, since $x_i^s=0$.
Hence, $(i,j)$ does not carry flow if $j$ receives all its flow from some source $s\neq i$.
Likewise, \eqref{eq:tMCCIFP} implies no flow along arc $(i,j)$ if $j\in T$ and $x_i^t=1$ for some sink $t\in T_i\setminus\{j\}$.

Redundant valid inequalities are derived by applying the idea behind constraints \eqref{eq:sMCCIFP}--\eqref{eq:tMCCIFP}
to sets of sources or sinks.
Consider an arc $(i,j) \in A$. If $j$ receives all its flow from some source $s\in S'\subseteq S$, i.e., if $\sum_{s\in S'}x_j^s=1$,
then there is no flow along $(i,j)$ unless also $\sum_{s\in S'}x_i^s=1$.
Therefore, the inequalities
\begin{eqnarray}
 & f_{ij}\leq h_{ij}\left(1+\sum_{s\in S'}(x_i^s-x_j^s)\right) & (i,j)\in A, S'\subseteq S \label{eq:visMCCIFP} \\
 & f_{ij}\leq h_{ij}\left(1+\sum_{t\in T'}(x_j^t-x_i^t)\right) & (i,j)\in A, T'\subseteq T \label{eq:vitMCCIFP}
\end{eqnarray}
are valid.
Alternatively, \eqref{eq:visMCCIFP}--\eqref{eq:vitMCCIFP} are derived from the observation that along any path $(i_1,\ldots,i_m)$ in $D(\boldsymbol{f})$,
$x_{i_k}^u$ is non-increasing with $k$ for all $u\in S$, and non-decreasing for all $u\in T$.
This follows directly from \eqref{eq:sMCCIFP}--\eqref{eq:tMCCIFP}.
The observation also implies validity of the combined inequality
\begin{equation}
	\displaystyle{f_{ij}\leq h_{ij}\left(1+\sum_{s\in S'}(x_i^s-x_j^s) + \sum_{t\in T'}(x_j^t-x_i^t)\right)}, ~(i,j)\in A, S'\subseteq S, T'\subseteq T, \label{eq:vistMCCIFP}
\end{equation}

\noindent
which is a generalization of \eqref{eq:visMCCIFP}--\eqref{eq:vitMCCIFP}.
The potential of inequality \eqref{eq:visMCCIFP} to strengthen \eqref{eq:sMCCIFP}, and \eqref{eq:vitMCCIFP} to strengthen \eqref{eq:tMCCIFP}, are addressed next:
\begin{proposition} \label{prop:strong}
Let $(\boldsymbol{f},\boldsymbol{x})$ be a feasible solution to the continuous relaxation of \eqref{eq:objMCCIFP}--\eqref{eq:xMCCIFP}.
If, for every arc $(i,j)\in A(\boldsymbol{f})$, it holds that $x_i^s\geq x_j^s$ for all $s\in S_j$ and $x_i^t\leq x_j^t$ for all $t\in T_i$,
then $\boldsymbol{f}$ is weakly confluent-influent.
\end{proposition}
\begin{proof}
The proof is by contradiction.
	Assume $\boldsymbol{f}$ is not weakly confluent-influent. Then, there exist a node $p\in P$,
	distinct sources $s_0,s_1\in S_p$, distinct sinks $t_0,t_1\in T_p$, and paths $Q_{k\ell}$ ($k,\ell=0,1$) in $D(\boldsymbol{f})$,
	where $Q_{k\ell}$ intersects $s_k$, $p$, and $t_{\ell}$.
	By definition, $x_{s_k}^{s_{1-k}}=x_{t_k}^{t_{1-k}}=0$ ($k=0,1$).
	As all arcs of $Q_{k\ell}$ are in $A(\boldsymbol{f})$, repeated application of $x_i^{s_1}\geq x_j^{s_1}$ to the arcs between $s_0$ and $p$ in $Q_{0\ell}$
	proves $x_p^{s_1}=0$. Analogous operations along $Q_{1\ell}$ yield $x_p^{s_0}=0$.
	Likewise, applying $x_i^{t_{1-\ell}}\leq x_j^{t_{1-\ell}}$ to the arcs between $t_{\ell}$ and $p$ in $Q_{k\ell}$ yields $x_p^{t_{\ell}}=0$ ($\ell=0,1$).
	Hence, by \eqref{eq:convMCCIFP}, $x_p^u>0$ for some $u\in (S_p\cup T_p)\setminus\{s_0,s_1,t_0,t_1\}$.
	If $u\in T_p$, applying $x_i^u\leq x_j^u$ to arcs $(i,j)$ along, e.g., $Q_{k0}$ from $p$ to $t_0$ yields $x_{t_0}^u>0$, which is a contradiction.
	An analogous analysis applies if $u\in S_p$.
\qed
\end{proof}

Consider an optimal solution $(\boldsymbol{f},\boldsymbol{x})$ to the continuous relaxation of \eqref{eq:objMCCIFP}--\eqref{eq:xMCCIFP},
and for all $(i,j)\in A$, let $\alpha_{ij}^s=x_i^s-x_j^s$ ($s\in S_j$) and $\alpha_{ij}^t=x_j^t-x_i^t$ ($t\in T_i$).
If $\boldsymbol{f}$ is weakly confluent-influent, then it is also an optimal flow in the MILP model.
Even if $\boldsymbol{x}$ is fractional, it can easily be adjusted to integer values consistent with $\boldsymbol{f}$.
Otherwise, Proposition \ref{prop:strong} shows that there exists some $(i,j)\in A(\boldsymbol{f})$ at which $\alpha_{ij}^u<0$ for some $u\in S_j\cup T_i$.
If there exist two distinct $u,v\in S_j\cup T_i$ such that $\alpha_{ij}^u,\alpha_{ij}^v<0$,
formulation \eqref{eq:objMCCIFP}--\eqref{eq:xMCCIFP} is strengthened by adding the inequality \eqref{eq:vistMCCIFP} corresponding to $S'\cup T'=\{u,v\}$.
If $f_{ij}>h_{ij}\left(1+\alpha_{ij}^u+\alpha_{ij}^v\right)$, then the added inequality cuts off $(\boldsymbol{f},\boldsymbol{x})$.

\subsection{Multi-commodity flow formulations} \label{sec:mciMILP}
For all $(i,j)\in A$ and $(s,t)\in S_i\times T_j$, denote by $f_{ij}^{st}$ the flow from source $s$ along arc $(i,j)$ to sink $t$,
and let $f_{ij}^{st}=0$ for $(s,t)\in (S\times T)\setminus(S_i\times T_j)$.
An upper bound, $h_{ij}^{st}$, on this flow is given as the maximum flow from $j$ to $i$ in the digraph $D_{ij}^{st}=(N,A_{ij}^{st})$,
where $A_{ij}^{st}=\left(A\cup\{(t,s)\}\right)\setminus\{(i,j)\}$, and the arc $(t,s)$ has capacity $h_{ts}=h_{ij}$.
% Define $b_u^{st}=b_u$ for $u\in\{s,t\}$, and $b_u^{st}=0$ for all $u\in (S\cup T)\setminus\{s,t\}$.
Let $\boldsymbol{f}^{st}=\left(f_{ij}^{st}\right)_{(i,j)\in A}$.
% $\boldsymbol{h}^{st}=\left(h_{ij}^{st}\right)_{(i,j)\in A}$, and
% $\boldsymbol{b}^{st}=\left(b_u^{st}\right)_{u\in S\cup T}$.

Imposing $f_{ij}^{st}=0$ unless $x_i^s+x_i^t=x_j^s+x_j^t=1$ and $x_i^t+x_j^s\leq 1$, a multi-commodity flow formulation for the weak MCCIFP is given as:

\begin{eqnarray}
\min & \sum_{(i,j)\in A}c_{ij}f_{ij} \label{eq:objmMCCIFP} \\
     & f_{ij} = \sum_{s\in S_i}\sum_{t\in T_j}f_{ij}^{st} & (i,j)\in A \label{eq:fmMCCIFP} \\
     & 0\leq f_{ij}^{st}\leq h_{ij}^{st}\left(x_i^s+x_i^t\right) & (i,j)\in A, s\in S_i, t\in T_j \label{eq:imMCCIFP} \\
     & 0\leq f_{ij}^{st}\leq h_{ij}^{st}\left(x_j^s+x_j^t\right) & (i,j)\in A, s\in S_i, t\in T_j \label{eq:jmMCCIFP} \\
     & 0\leq f_{ij}^{st}\leq h_{ij}^{st}\left(2-x_i^t-x_j^s\right) & (i,j)\in A, s\in S_i, t\in T_j \label{eq:ijmMCCIFP} \\
     & \eqref{eq:convMCCIFP}, \eqref{eq:fMCCIFP}-\eqref{eq:xMCCIFP}\label{eq:mMCCIFP}
\end{eqnarray}

\subsection*{Acknowledgment}
\small{The work of the second author has been funded by grants from the Spanish state (TIN2012-37483-C03-01 and TIN2015-66680), Junta de Andaluc\'{\i}a (P11-TIC-7176),
in part financed by the European Regional Development Fund (ERDF).}

\bibliographystyle{elsarticle-num}
\bibliography{ref.bib}

\end{document}
